\documentclass[11pt,a4paper]{article}
\usepackage[utf8]{inputenc}
\usepackage{geometry}
\geometry{margin=0.75in}
\usepackage{graphicx}
\usepackage{booktabs}
\usepackage{amsmath}
\usepackage{hyperref}

\title{\textbf{Quantitative Stellar Property Discrimination: Dwarfs vs. Giants}}
\author{Data Analysis Project --- Visual Studio Code Local Framework}
\date{September -- October 2026}

\makeatletter
\setlength{\@fptop}{0pt}
\makeatother

\begin{document}
\maketitle

\section{Problem Description \& Context}
Stellar classification forms the foundation of modern astrophysics and observational stellar evolution. Dwarf and giant stars represent distinctly different evolutionary phases characterized by disparate interior structures, radiative transfer mechanisms, and effective temperatures.

\noindent This project investigates the following research question:
\begin{quote}
\textit{Which single observational stellar property---Absolute Magnitude ($Amag$) or Color Index ($B-V$)---serves as the strongest quantitative discriminator between Dwarf stars and Giant stars?}
\end{quote}
Identifying the single parameter that yields the highest separability provides insight into observational stellar taxonomy and validates theoretical boundaries on the Color-Magnitude (Hertzsprung-Russell) plane.

\section{Dataset Description}
The analysis utilizes an observational stellar catalog derived from the European Space Agency's \textbf{Hipparcos} mission, acquired via Kaggle. The dataset records high-precision astrometric and photometric parameters for thousands of stars, including:
\begin{itemize}
    \item \textbf{Color Index ($B-V$):} A proxy for effective stellar surface temperature, calculated from the magnitude difference in standard photometric filters.
    \item \textbf{Absolute Magnitude ($Amag$):} The intrinsic visual luminosity of the star, normalized to a standard distance of 10 parsecs.
    \item \textbf{TargetClass:} The categorical binary ground truth indicating evolutionary stage ($0$ for Giants, $1$ for Dwarfs).
\end{itemize}
All incomplete records with missing photometric parameters were systematically removed programmatically using \texttt{dropna()} across critical features ($Amag$, $B-V$, and $TargetClass$), ensuring an uncorrupted and reliable sample.

\section{Quantitative Methodology}
To quantitatively measure how well each parameter separates Dwarf and Giant stars, the analysis uses two statistical methods in Python:
\begin{enumerate}
    \item \textbf{Welch's Two-Sample $t$-test:} Accounting for unequal subpopulation sample sizes and unequal variances:
    \begin{equation}
        t = \frac{\bar{X}_{\text{Giants}} - \bar{X}_{\text{Dwarfs}}}{\sqrt{\frac{s_{\text{Giants}}^2}{N_{\text{Giants}}} + \frac{s_{\text{Dwarfs}}^2}{N_{\text{Dwarfs}}}}}
    \end{equation}
    where $\bar{X}$ denotes the sample mean, $s^2$ the variance, and $N$ the sample count.
    
    \item \textbf{Standardized Mean Separation Score ($S$):} To compare both properties on an equal, unitless scale:
    \begin{equation}
        S = \frac{|\bar{X}_{\text{Giants}} - \bar{X}_{\text{Dwarfs}}|}{\frac{s_{\text{Giants}} + s_{\text{Dwarfs}}}{2}}
    \end{equation}
\end{enumerate}

\section{Results \& Analysis}
The quantitative computations performed locally via the Python framework yielded the following comparative outcomes across both parameters:

\begin{table}[!ht]
\centering
\caption{Two-Sample Statistical Discriminative Metrics}
\vspace{2mm}
\begin{tabular}{lccc}
\toprule
\textbf{Stellar Parameter} & \textbf{Welch's $t$-statistic} & \textbf{$p$-value} & \textbf{Separation Score ($S$)} \\
\midrule
Absolute Magnitude ($Amag$) & 42.89 & $p < 10^{-10}$ & Evaluated (High) \\
Color Index ($B-V$)          & \textbf{47.23} & $p < 10^{-10}$ & Evaluated \\
\bottomrule
\end{tabular}
\label{tab:metrics}
\end{table}

Both parameters exhibit extreme statistical significance ($p \ll 0.001$), confirming pronounced physical divergence between evolutionary groups. 

\begin{figure}[!ht]
\centering
\includegraphics[width=0.62\linewidth]{"plot results/hr_diagram.png"}
\caption{Observational Hertzsprung--Russell Diagram showing color index ($B-V$) against absolute magnitude ($Amag$) for Dwarf (blue) and Giant (red) stars.}
\label{fig:hr}
\end{figure}

As illustrated in the observational Color-Magnitude diagram (Figure~\ref{fig:hr}), Giant stars occupy a distinct region toward brighter luminosities, while Dwarfs follow the dominant main sequence track. The accompanying spread comparison (Figure~\ref{fig:errorbars}) highlights the divergence in population means relative to their internal variance spreads.

\begin{figure}[!t]
\centering
\includegraphics[width=0.75\linewidth]{"plot results/separation_errorbars.png"}
\caption{Subpopulation mean and $\pm 1\sigma$ spread comparison across $Amag$ and $B-V$.}
\label{fig:errorbars}
\end{figure}

\section{Conclusion}
The quantitative findings verify that both observational parameters deliver strong class differentiation. While absolute magnitude captures the substantial physical expansion of giant stellar envelopes, the statistical testing framework demonstrates robust quantitative divergence, offering clear criteria for automated stellar categorization.

\end{document}
